\documentclass[11pt]{article}
\usepackage[margin=25mm]{geometry}
\usepackage{fontspec}
\setmainfont{Latin Modern Roman}
\usepackage{amsmath,amssymb,mathtools}
\usepackage{unicode-math}
\setmathfont{Latin Modern Math}
\usepackage{xurl,hyperref}
\hypersetup{colorlinks=true,urlcolor=blue}
\setcounter{secnumdepth}{0}
\setlength{\parindent}{0pt}
\setlength{\parskip}{5pt}
\setlength{\emergencystretch}{3em}
\newcommand{\sref}[2]{\hyperlink{source:#1:#2}{[S#2]}}
\newcommand{\cataloguescope}[1]{\paragraph{Recorded target.}#1}
\begin{document}
\section{Kirby Problem 4.120}\label{statement}

\subsection{Definitions and mathematical statement}

\paragraph{Shared notation.}$\mathbb N=\{1,2,\ldots\}$, and $\mathbb Z,\mathbb Q,\mathbb R,\mathbb C$ denote the integers, rationals, reals and complex numbers. Any other conventions are specified locally.


An achiral Lefschetz pencil on a smooth closed four-manifold $X$ consists of a finite base set $B\subset X$ and a map $f:X\setminus B\to S^2$ with finitely many isolated nodal critical points. Up to smooth coordinate changes, its local critical-point model is $(z_1,z_2)\mapsto z_1^2+z_2^2$, and at a base point its local model is $[z_1:z_2]\in\mathbb{CP}^1$. Coordinates in these models may reverse orientations, in contrast to a chiral pencil. Away from these points the map is a submersion. The source allows no critical points and/or no base points, and allows nonorientable $X$.
\paragraph{Statement recorded in the supplied source.}
Characterize the closed smooth four-manifolds admitting an achiral Lefschetz pencil. In particular, does every closed simply connected smooth four-manifold admit such a pencil on the original manifold itself? \sref{statement}{2}
\subsection{Short English statement}
Determine which closed smooth four-manifolds support an achiral Lefschetz pencil, especially whether every simply connected one does.
\subsection{Sources}
\begin{description}
\item[\hypertarget{source:statement:1}{S1}] ULAM.AI and the UnsolvedMath Contributors, \emph{UnsolvedMath: A Curated Collection of Open Mathematics Problems}, catalogue entry. CC BY 4.0. \url{https://huggingface.co/datasets/ulamai/UnsolvedMath}.
\item[\hypertarget{source:statement:2}{S2}] R. İnanç Baykur, Robion C. Kirby and Daniel Ruberman (eds.), \emph{K3: A New Problem List in Low-Dimensional Topology}, Mathematical Surveys and Monographs 295, American Mathematical Society (2026), Problem 4.120, author preliminary version. \url{https://math.berkeley.edu/sites/default/files/surv-295-ruberman-watermarked-author-pdf.pdf}.
\end{description}
\label{end:statement}
\end{document}
